% =====================================================================
%  04  熱統合 — CC（与熱・受熱複合線）と GCC を横並び。
%   ヘッダー帯に「CC」「GCC」を表示し、図内のタイトルは削除して図を最大化。
%   図は帯直下〜スライド最下部いっぱいに配置（x軸ラベルを下端へ密着）。
%  source_memo: 09_heat_integration.tex / 10_optimization.tex §経済性。
%   図 figures/make_composite_slide.py（gcc_bestpoint.json 由来）。
% =====================================================================
{%  この frame だけヘッダー帯を CC / GCC に差し替える（group でスコープ）
\setbeamertemplate{headline}{%
  \ifnum\insertframenumber>0\relax
    \SetBarColor
    \begin{beamercolorbox}[wd=\paperwidth, ht=4.5ex, dp=1.8ex,
                           leftskip=1.2em, rightskip=1.2em]{headerbar}%
      \color{white}\bfseries\large
      CC・GCC%
      \hfill
      \insertframenumber/\MainTotal%
    \end{beamercolorbox}%
    \vskip-2.6ex
  \fi
}%
\begin{frame}[t]
  \vspace*{0.1ex}
  % 図を縦に収めて x 軸ラベルの見切れを防ぐため、全幅から少し縮めて中央配置。
  \centering
  \includegraphics[width=\dimexpr\linewidth-2mm\relax]{composite_slide}
\end{frame}
}
